\documentclass[a4paper,11pt]{article}

% -----------------------------------------------------------------------
% Packages
% -----------------------------------------------------------------------
\usepackage{amsmath,amssymb,amsthm}
\usepackage{hyperref}
\usepackage{booktabs}
\usepackage{microtype}
\usepackage{xcolor}
\usepackage{listings}

% -----------------------------------------------------------------------
% Theorem environments
% -----------------------------------------------------------------------
\newtheorem{theorem}{Theorem}[section]
\newtheorem{lemma}[theorem]{Lemma}
\newtheorem{corollary}[theorem]{Corollary}
\newtheorem{definition}[theorem]{Definition}
\newtheorem{proposition}[theorem]{Proposition}
\newtheorem{remark}[theorem]{Remark}

% -----------------------------------------------------------------------
% Title
% -----------------------------------------------------------------------
\title{%
  Structural Admissibility for Verification Systems:\
  Observation, Warrant Debt, and Factorisation Repair%
}

\author{%
  % TODO: author names
}

\date{}

% -----------------------------------------------------------------------
\begin{document}
\maketitle

% -----------------------------------------------------------------------
\begin{abstract}
Safety claims are usually made through an observation interface: a monitor, a
log, a trace, a sensor, or an abstraction. Such interfaces do not preserve every
concrete distinction. This matters when a safety predicate distinguishes states
that the interface identifies. We call a predicate admissible relative to an
observation map when it factors through that map. If factorisation fails, the
missing distinction cannot be recovered by post-processing the observation. We
formalise this criterion in Rocq. The development contains 39 compiled files, no
admitted lemmas, and eight named axioms. The case studies are deliberately
small. They cover consensus under partition, GPU-style relaxed memory
observations, and physical sensor boundary models. They do not claim to be full
models of those domains. Their purpose is narrower: to show how the same
factorisation failure can appear in discrete, concurrent, and physical settings.
The remedy is also ordinary. Safety practice already uses both restriction and
refinement: either weaken or condition the claim so that the available
observation supports it, or improve the observation so that the relevant
distinction is visible. The contribution is to state this discipline explicitly,
and to show why treating such failures as mere noise is unsafe.
\end{abstract}

% -----------------------------------------------------------------------
\section{Introduction}
\label{sec:introduction}
% -----------------------------------------------------------------------

Verification rarely sees the system directly. It sees an abstraction, a trace, a
log, a monitor state, a sensor reading, or a reconstructed state estimate. Each
is an observation map. Each preserves some distinctions and collapses others.
This is not an error. It is the normal condition under which verification and
runtime assurance operate.

The risk begins when a safety predicate depends on a distinction that the
observation does not preserve. Then the safety claim is stronger than the
interface through which it is made. In routine engineering this problem is often
handled pragmatically. One restricts the operating envelope, adds a sensor,
changes a monitor, rejects an over-strong claim, or records an extra field in a
log. These are familiar moves. The paper gives them a simple structural test.

A safety predicate is admissible relative to an observation map exactly when it
factors through that map. Equivalently, states that have the same observation
must agree on the predicate. If two concrete states look the same to the
observation but differ with respect to the safety predicate, the claim is not
supported by that observation. No downstream classifier, dashboard, optimiser,
or monitor can repair the loss. The information has already been collapsed.

This point is modest, but important. Many admissibility failures will look like
noise at the observational level. A sensor appears imprecise. A trace appears
partial. A local execution appears harmless. A monitor appears inconclusive. The
paper argues that this is sometimes the wrong diagnosis. The problem may not be
noise. It may be a failed factorisation.

The paper does not propose a new general-purpose verification algorithm. It
does not claim that all safety reasoning reduces to factorisation. It isolates a
necessary condition on safety claims made through observation interfaces. When
that condition fails, the claim must be restricted, rejected, or the observation
must be refined.

\paragraph{Contributions.}

\begin{itemize}
\item We define structural admissibility as factorisation of a safety predicate
through an observation map.
\item We prove the non-recoverability principle: once an observation collapses a
safety-relevant distinction, post-processing cannot separate it.
\item We mechanise the framework in Rocq, with 39 compiled source files, no
admitted lemmas, and eight named axioms.
\item We give small case studies in consensus, GPU-style memory observations,
and physical sensor models.
\item We distinguish two repair forms: predicate rejection or restriction, and
observational refinement. These are formally parallel but not operationally
symmetric.
\end{itemize}

Section~\ref{sec:setting} states the factorisation criterion. Section~\ref{sec:warrant}
uses it to define warrant debt and non-recoverability. Section~\ref{sec:recovery}
explains restriction and refinement. Section~\ref{sec:rocq} describes the Rocq
development and its assumptions. Section~\ref{sec:casestudies} reports the case
studies. Sections~\ref{sec:related} to~\ref{sec:conclusion} place the result,
state its limits, and draw the practical lesson.

% -----------------------------------------------------------------------
\section{Mathematical Setting}
\label{sec:setting}
% -----------------------------------------------------------------------

Let \(X\) be a concrete state space, \(O\) an observation space, and
\(M : X \to O\) an observation map. Let \(\Phi : X \to \mathsf{Prop}\) be a
safety predicate.

\subsection{Observation-induced equivalence}

The observation map \(M\) induces an equivalence on concrete states:
\[
  x \sim_M y \quad\text{iff}\quad M(x) = M(y).
\]
States that agree under \(M\) are indistinguishable from the perspective of
the observation interface.

\subsection{Admissibility}

\begin{definition}[Admissibility]
\label{def:admissible}
The predicate \(\Phi\) is \emph{admissible} relative to \(M\) if and only if
\[
  M(x) = M(y) \Longrightarrow \bigl(\Phi(x) \leftrightarrow \Phi(y)\bigr).
\]
\end{definition}

Equivalently, \(\Phi\) is constant on the fibres of \(M\).

\subsection{Factorisation}

\begin{definition}[Factorisation]
\label{def:factorisation}
The predicate \(\Phi\) \emph{factors through} \(M\) if there exists
\(\widehat{\Phi} : O \to \mathsf{Prop}\) such that
\[
  \forall x : X,\quad \Phi(x) \leftrightarrow \widehat{\Phi}(M(x)).
\]
\end{definition}

\subsection{The main equivalence}

\begin{theorem}[Admissibility iff factorisation]
\label{thm:backbone}
The predicate \(\Phi\) is admissible relative to \(M\) if and only if it
factors through \(M\).
\end{theorem}

\begin{proof}
This is the standard quotient argument. If \(\Phi\) factors through \(M\), then
\(M(x)=M(y)\) immediately implies that \(\Phi(x)\) and \(\Phi(y)\) have the same
truth value. Conversely, if \(\Phi\) is constant on the fibres of \(M\), define
\(\widehat{\Phi}\) on each observable value by choosing any representative in its
fibre. The constancy condition makes this definition independent of the chosen
representative. In the Rocq development this is the theorem
\texttt{admissible\_iff\_constant\_on\_kernel}.
\end{proof}

This equivalence is the formal spine of the paper. Later results use it; they do
not replace it.

% -----------------------------------------------------------------------
\section{Warrant Debt and Non-Recoverability}
\label{sec:warrant}
% -----------------------------------------------------------------------

\subsection{Warrant debt}

\begin{definition}[Warrant debt]
\label{def:debt}
An observation map \(M\) has \emph{warrant debt} with respect to predicate
\(\Phi\) when \(\Phi\) is not admissible relative to \(M\). Concretely, there
exist \(x, y : X\) such that
\[
  M(x) = M(y),\quad \Phi(x),\quad \neg\,\Phi(y).
\]
\end{definition}

Warrant debt is the gap between what the observation supports and what the
predicate requires. The term is meant literally. A safety claim made through
\(M\) borrows distinctions that \(M\) does not provide.

\subsection{Non-recoverability}

\begin{theorem}[Non-recoverability]
\label{thm:no-recovery}
For any map \(f : O \to Y\),
\[
  M(x) = M(y) \Longrightarrow f(M(x)) = f(M(y)).
\]
Consequently, if \(M\) collapses \(x\) and \(y\), no downstream map from
observations can separate them.
\end{theorem}

\begin{proof}
By congruence of equality under \(f\). In the Rocq development this is the
non-recoverability lemma for post-processing maps over observations.
\end{proof}

\begin{remark}
No dashboard, classifier, monitor, optimiser, or post-processor can recover a
distinction that the observation map has already destroyed. The information is
not latent. It is absent.
\end{remark}

\subsection{Why this can be mistaken for noise}

At the observational level, an inadmissibility witness may look unremarkable.
Two executions have the same local trace. Two sampled trajectories have the
same visible readings. Two distributed states have the same local view. It is
tempting to describe the disagreement as measurement noise, modelling
imprecision, or a benign abstraction artefact. Sometimes that is correct. But
when the predicate separates two states that the observation identifies, the
problem is structural. The observation is not merely noisy with respect to the
claim. It is insufficient for the claim.

This is the practical reason for naming warrant debt. The name does not add a
new proof principle. It marks a common mistake: treating a failed observation
boundary as if it were only an implementation nuisance.

\subsection{Consequence for verification}

If a proof of \(\Phi\) requires a distinction invisible to \(M\), the proof is
not a proof through \(M\). It is a proof through a stronger, implicit
observation. The framework makes this precise. An inadmissible predicate cannot
be made admissible by adding an assumption that bypasses the observation
boundary. The assumption merely moves the missing observation into the proof.

% -----------------------------------------------------------------------
\section{Recovery: Restrict or Refine}
\label{sec:recovery}
% -----------------------------------------------------------------------

At the level of the factorisation criterion, there are two principled repair
forms.

\subsection{Predicate rejection or restriction}

One may replace \(\Phi\) by a predicate \(\Psi\) that is admissible relative to
the existing observation \(M\). Depending on the application, this may mean
weakening the guarantee, restricting the class of states under consideration,
strengthening the operating envelope, using robust margins, or limiting the
class of executions to those whose structure makes the observation adequate.

Rejection does not always mean weakening the predicate in a crude sense. It
means refusing to assert a claim whose truth is not determined by the available
observation. Restriction is the constructive version of that refusal.

\subsection{Observational refinement}

Alternatively, one may replace \(M\) by a richer observation \(M' : X \to O'\)
together with a projection \(\pi : O' \to O\) such that
\[
  M = \pi \circ M'.
\]
The refined observation \(M'\) preserves the old observation while exposing
additional structure. If \(\Phi\) factors through \(M'\), admissibility has been
recovered by changing the observation boundary rather than the predicate.

\subsection{Formal duality is not operational symmetry}
\label{subsec:not-symmetric}

Restriction and refinement are formally parallel repair forms. They are not the
same engineering act. Restriction keeps the observation map fixed and changes
the claim. Refinement keeps the claim fixed and changes the observation. Both
can restore factorisation. They pay different costs.

Restriction pays in guarantee strength. The resulting claim may be narrower,
conditional, local, robust only within a margin, or valid only under a specified
synchronisation discipline. Refinement pays in observational burden. The system
must expose more structure, log more information, add a sensor, preserve a
trace, or reconstruct a richer state.

This distinction addresses an asymmetry in the case studies. The consensus and
GPU examples primarily mechanise refinement. The physical-systems example
mechanises both robust restriction and refinement. This is not presented as a
failure of the framework. It is part of the lesson. The repair that is natural in
one domain may be artificial in another. A GPU memory example can often be made
admissible by moving from thread-local observation to a global execution
observation. A physical sensor example may instead require explicit margins,
boundary constraints, and a richer account of connection. The same formal
failure is present, but the cost of repair is domain-specific.

\subsection{Everyday safety procedures already use both repairs}
\label{subsec:safety-practice}

The distinction is not exotic. Everyday safety procedures already use both
moves. A restricted operating envelope is predicate restriction: the claim is no
longer that all states are safe, but that states inside a stated envelope are
safe. An added sensor is observational refinement: the claim is kept, but the
interface is strengthened. A lockout condition restricts admissible operation. A
new alarm channel refines observation. A conservative threshold restricts the
claim to a margin. A higher resolution measurement refines the observation.

The paper's point is not that engineers do not know this. They do. The point is
that verification texts and monitor designs can obscure the distinction. Once
the distinction is obscured, an unsupported safety predicate can be smuggled in
as a harmless assumption, a modelling convention, or noise. The factorisation
criterion is a discipline against that mistake.

\subsection{No invisible axioms}

\begin{quote}
When a predicate fails to factor through the current observation map, the proof
must not be rescued by an assumption that bypasses the observation boundary. The
claim must be restricted to what the observation supports, or the observation
must be refined until the missing distinction is exposed.
\end{quote}

The case studies in Section~\ref{sec:casestudies} illustrate this discipline.
They should not be read as claiming that restriction and refinement are equally
natural in every domain.

% -----------------------------------------------------------------------
\section{Rocq Development}
\label{sec:rocq}
% -----------------------------------------------------------------------

\subsection{Library structure}

The mechanisation is organised as follows:

\begin{description}
\item[\texttt{Foundation/}] Relations, orders, lattices, quotients, topology,
mereotopology, and metric resolution.
\item[\texttt{Admissibility/}] Observation maps, the factorisation criterion,
kernel pairs, refinement, collapse witnesses, warrant debt, and horizon debt.
\item[\texttt{Runtime/}] Monitor states, structural fatigue, metastable closure,
rupture certificates, recovery, and runtime soundness.
\item[\texttt{Automation/}] Custom tactics for admissibility, refinement,
collapse, and certificate reasoning.
\item[\texttt{CaseStudies/Consensus/}] Network model, partition collapse,
consensus admissibility.
\item[\texttt{CaseStudies/GPU/}] Memory model, GPU admissibility.
\item[\texttt{CaseStudies/PhysicalSystems/}] Boundary conditions, thermal model,
sensor collapse, PDE admissibility.
\item[\texttt{Extraction/}] Extractable types, monitor extraction, certificate
extraction, extraction correctness.
\end{description}

\subsection{Formal status}

The accompanying Rocq development consists of 39 source files, all of which
compile to \texttt{.vo} files. It contains no admitted lemmas. The remaining
logical strength of the development is confined to eight named axioms. They are
not hidden. They are part of the mathematical specification.

\subsection{Axioms and their role}
\label{subsec:axioms-role}

The eight axioms are background principles, not incomplete proofs. Three concern
primitive connection in mereotopology:

\begin{description}
\item[\texttt{C\_refl}] Every region is connected to itself.
\item[\texttt{C\_sym}] Connection is symmetric.
\item[\texttt{ntpp\_irreflexive}] No region is a non-tangential proper part of
itself.
\end{description}

The axiom \texttt{ntpp\_irreflexive} is an RCC-style mereotopological
assumption. It is not derivable from \texttt{C\_refl} and \texttt{C\_sym}
alone. It is therefore stated explicitly. The derived result
\texttt{ntpp\_antisym} is proved from it.

The remaining five axioms are application-level ordering and coherence
constraints for the physical-systems case study and runtime monitor parameters.
They should be listed in the final submitted version by their exact Rocq names.
The point of the list is not to increase apparent rigour. It is to let the
reader see where the physical model obtains its strength.

\begin{table}[h]
\centering
\begin{tabular}{lll}
\toprule
Axiom class & Role in the model & Why it is explicit \\
\midrule
Connection reflexivity & Regions connect to themselves & Base mereotopology \\
Connection symmetry & Connection has no direction & Base mereotopology \\
NTPP irreflexivity & Proper interior containment is not self-containment & Boundary discipline \\
Physical ordering constraint & Prevents invalid ordering of physical states & Model coherence \\
Physical boundary constraint & Keeps boundary reasoning non-degenerate & Sensor admissibility \\
Monitor threshold constraint & Orders runtime thresholds coherently & Runtime soundness \\
Recovery margin constraint & Separates robust restriction from unsafe closure & Repair soundness \\
Horizon or timing constraint & Connects admissibility to intervention time & Runtime interpretation \\
\bottomrule
\end{tabular}
\caption{Axiom roles. Replace the last five descriptive labels by the exact Rocq axiom names from \texttt{ASSUMPTIONS.md} before submission.}
\label{tab:axiom-roles}
\end{table}

This table answers a likely objection. The continuous case does more
mathematical work than the discrete examples. It uses boundary assumptions and
ordering constraints that have no direct analogue in a thread-local trace or a
partitioned local view. The common claim is not that the domains have the same
mathematics. They do not. The common claim is that each contains an observation
map, a predicate, and a possible failure of factorisation.

\subsection{Mechanisation scope}

The development mechanises the factorisation criterion and its equivalence with
admissibility, the non-recoverability principle, the runtime monitor theory
including rupture certificates, and the three case-study examples. It does not
claim to be a full mechanisation of Paxos, a production GPU memory model, or a
continuous control-system specification.

% -----------------------------------------------------------------------
\section{Case Studies}
\label{sec:casestudies}
% -----------------------------------------------------------------------

% case-studies.tex
% Integrated case-study section for the structural admissibility paper.
% Insert into the main file with:
%   % case-studies.tex
% Integrated case-study section for the structural admissibility paper.
% Insert into the main file with:
%   % case-studies.tex
% Integrated case-study section for the structural admissibility paper.
% Insert into the main file with:
%   \input{paper/case-studies}
% at the point where the case-study section begins.

\section{Case Studies: Refine or Reject}

The development contains three case-study families: consensus under network
partition, GPU-style relaxed memory, and cyber-physical sensor models. These
examples are deliberately drawn from different engineering domains. Their common
structure is the point. In each case, a safety predicate is naturally formulated
over a concrete state space~\(X\), while the available observation map
\(M : X \to O\) exposes only a coarser interface. The case study then asks
whether the predicate factors through that observation.

The answer is negative for the initial coarse observation in each family. The
coarse observation identifies concrete states that disagree on the relevant
safety predicate. The admissibility failure is therefore structural rather than
accidental. Recovery takes one of two forms: replace the predicate by one that
is meaningful at the observational grain, or refine the observation so that the
safety-relevant distinction is no longer collapsed.

\begin{table}[t]
\centering
\caption{Case-study structure: factorisation failure and admissibility recovery.}
\label{tab:case-study-factorisation}
\small
\begin{tabular}{p{1.8cm}p{2.8cm}p{3.5cm}p{2.9cm}}
\hline
Case study & Hidden distinction & Coarse failure & Recovery \\
\hline
Consensus
  & Partition and quorum structure
  & Local observations do not determine global quorum safety
  & Restrict to local quorum belief; or enrich with quorum-safety data \\[4pt]
GPU memory
  & Buffer, cache, and ordering structure
  & Value traces do not determine consistency behaviour
  & Restrict to DRF executions; or enrich with global ordering metadata \\[4pt]
Physical systems
  & Inter-sample and boundary behaviour
  & Sampled values do not determine continuous safety
  & Use robust observational margins; or enrich sensing and boundary data \\
\hline
\end{tabular}
\end{table}

\paragraph{Mechanisation status.}
The three case-study families instantiate the same factorisation pattern.
The consensus case study mechanises both repair paths: the refinement
path restores admissibility for \(\texttt{quorum\_safe}\) by enriching
the observation (\texttt{ConsensusAdmissibility.v}), and the restriction
path yields an admissible local-quorum-belief predicate at the coarse
observational grain (\texttt{ConsensusRestriction.v}). The GPU case study
mechanises both repair paths: refinement by enriched global execution
observation (\texttt{GPUAdmissibility.v}), and restriction to the subtype
of data-race-free executions (\texttt{GPURestriction.v}). The
physical-systems case study mechanises both paths: a restricted predicate
is admissible at the coarse sensor grain, and an enriched sensor observation
restores admissibility for the stronger predicate. The development
mechanises predicate restriction in all three domains and observational
refinement in all three.

% -----------------------------------------------------------------------
\subsection{Consensus Under Partition}

The consensus case study concerns quorum safety under network partition. The
concrete state records enough information to determine whether the relevant
quorum-safety predicate holds. A coarse observation, however, may expose only
local status information. Such an observation does not determine global quorum
safety. Two concrete states may agree on the local observation while disagreeing
on whether a quorum is safely established.

This failure is not a defect in the proof script. It is the phenomenon the
framework is intended to detect. A uniformity principle for reachable alive
nodes fails in the partition model: different nodes may have different
reachable alive counts. The admissibility proof cannot be closed by that
route.

The mechanised consensus case study proves both repair paths.

\textit{Refinement.} The observation record \texttt{ConsensusObs} contains the
field \texttt{co\_quorum\_safe}, computed by \texttt{consensus\_obs} from the
underlying consensus state. The proof of admissibility is then immediate:
equality of observations entails equality of the quorum-safety component by
projection. The safety predicate becomes admissible because the observation
carries the safety-relevant distinction directly.

\textit{Restriction.} Rather than asserting global quorum safety, one can
instead assert only what the local observation supports: the node's own quorum
belief, \texttt{local\_quorum\_belief~n}, a projection of \texttt{lv\_quorum}
from the local view record. This predicate is admissible relative to
\texttt{local\_node\_obs~n} because it is a projection of the observation.
The proof requires no axiom. The theorem \texttt{restriction\_is\_genuinely\_weaker}
confirms that local belief and global quorum safety can disagree, so this is a
genuine weakening and not a reformulation.

The theorem \texttt{both\_repairs\_are\_available} states the formal duality
in a single proposition: both \texttt{local\_quorum\_belief~n} and
\texttt{quorum\_safe} are admissible, the former relative to the coarse local
observation, the latter relative to the enriched consensus observation.

% -----------------------------------------------------------------------
\subsection{GPU-Style Memory Observations}

The GPU-style case study concerns the gap between value observations and memory
consistency. A coarse observation may record only the values read and written by
threads. The concrete execution state, however, contains additional ordering and
visibility structure: buffering, cache state, scope, and synchronisation
metadata. A predicate such as sequential consistency of the concurrent execution
is not, in general, determined by the value trace alone.

The admissibility failure has the same form as in the consensus case. Two
executions may have the same observed value trace for a given thread while
disagreeing on whether the execution is sequentially consistent. The observation
map has collapsed a distinction on which the safety predicate depends. The
predicate therefore fails to factor through the thread-local observation.

The GPU case study mechanises both repair paths. Refinement restores
admissibility by enriching the observation with global execution structure.
Restriction keeps the thread-local observation fixed but restricts the domain
to data-race-free executions. On that restricted domain, the
sequential-consistency predicate is stable across observation fibres,
conditional on the DRF-SC bridge stated in \texttt{GPURestriction.v}.

% -----------------------------------------------------------------------
\subsection{Physical Sensing and Boundary Conditions}

The physical-systems case study concerns safety predicates over continuous or
boundary-sensitive states. A controller or monitor rarely observes the full
physical trajectory. It receives sampled, discretised, or otherwise compressed
sensor data. A predicate may require that a variable remain within a safe region
throughout an interval, while the observation records only values at selected
sampling points or coarse boundary summaries.

The resulting admissibility failure is again a failure of factorisation. Two
physical trajectories may induce the same sampled observation while disagreeing
on the safety predicate. One trajectory may remain inside the safe envelope,
while another may violate the envelope between samples or across an unresolved
boundary. The sampled observation is therefore insufficient to determine the
predicate.

This case study mechanises both recovery paths. The coarse sensor observation is
proved inadmissible for the original safety predicate. A restricted predicate
meaningful at the coarse observational grain is then shown to be admissible
relative to that observation. Finally, an enriched sensor observation that adds
the missing physical data restores admissibility for the stronger predicate. The
physical-systems case study therefore demonstrates the full design choice: one
may weaken the guarantee to match the observation, or strengthen the observation
to match the guarantee.

% -----------------------------------------------------------------------
\subsection{Common Structure}

The three case-study families differ in engineering content but not in logical
shape. Consensus hides partition and quorum structure. GPU-style memory hides
ordering and visibility structure. Physical sensing hides inter-sample or
boundary behaviour. In each case, the initial observation map is too coarse for
the safety predicate. The predicate fails to factor because the observation
identifies concrete states that the predicate separates.

This is the common admissibility defect. It is independent of whether the hidden
structure is distributed, concurrent, or physical. The repair is also common:
either choose a predicate that is meaningful at the given observational grain, or
refine the observation until the predicate becomes invariant on observation
fibres. The case studies therefore support the main methodological claim of the
paper. Admissibility is not merely a property of a predicate or a system. It is a
property of the relation between a predicate and an observation interface.

% at the point where the case-study section begins.

\section{Case Studies: Refine or Reject}

The development contains three case-study families: consensus under network
partition, GPU-style relaxed memory, and cyber-physical sensor models. These
examples are deliberately drawn from different engineering domains. Their common
structure is the point. In each case, a safety predicate is naturally formulated
over a concrete state space~\(X\), while the available observation map
\(M : X \to O\) exposes only a coarser interface. The case study then asks
whether the predicate factors through that observation.

The answer is negative for the initial coarse observation in each family. The
coarse observation identifies concrete states that disagree on the relevant
safety predicate. The admissibility failure is therefore structural rather than
accidental. Recovery takes one of two forms: replace the predicate by one that
is meaningful at the observational grain, or refine the observation so that the
safety-relevant distinction is no longer collapsed.

\begin{table}[t]
\centering
\caption{Case-study structure: factorisation failure and admissibility recovery.}
\label{tab:case-study-factorisation}
\small
\begin{tabular}{p{1.8cm}p{2.8cm}p{3.5cm}p{2.9cm}}
\hline
Case study & Hidden distinction & Coarse failure & Recovery \\
\hline
Consensus
  & Partition and quorum structure
  & Local observations do not determine global quorum safety
  & Restrict to local quorum belief; or enrich with quorum-safety data \\[4pt]
GPU memory
  & Buffer, cache, and ordering structure
  & Value traces do not determine consistency behaviour
  & Restrict to DRF executions; or enrich with global ordering metadata \\[4pt]
Physical systems
  & Inter-sample and boundary behaviour
  & Sampled values do not determine continuous safety
  & Use robust observational margins; or enrich sensing and boundary data \\
\hline
\end{tabular}
\end{table}

\paragraph{Mechanisation status.}
The three case-study families instantiate the same factorisation pattern.
The consensus case study mechanises both repair paths: the refinement
path restores admissibility for \(\texttt{quorum\_safe}\) by enriching
the observation (\texttt{ConsensusAdmissibility.v}), and the restriction
path yields an admissible local-quorum-belief predicate at the coarse
observational grain (\texttt{ConsensusRestriction.v}). The GPU case study
mechanises both repair paths: refinement by enriched global execution
observation (\texttt{GPUAdmissibility.v}), and restriction to the subtype
of data-race-free executions (\texttt{GPURestriction.v}). The
physical-systems case study mechanises both paths: a restricted predicate
is admissible at the coarse sensor grain, and an enriched sensor observation
restores admissibility for the stronger predicate. The development
mechanises predicate restriction in all three domains and observational
refinement in all three.

% -----------------------------------------------------------------------
\subsection{Consensus Under Partition}

The consensus case study concerns quorum safety under network partition. The
concrete state records enough information to determine whether the relevant
quorum-safety predicate holds. A coarse observation, however, may expose only
local status information. Such an observation does not determine global quorum
safety. Two concrete states may agree on the local observation while disagreeing
on whether a quorum is safely established.

This failure is not a defect in the proof script. It is the phenomenon the
framework is intended to detect. A uniformity principle for reachable alive
nodes fails in the partition model: different nodes may have different
reachable alive counts. The admissibility proof cannot be closed by that
route.

The mechanised consensus case study proves both repair paths.

\textit{Refinement.} The observation record \texttt{ConsensusObs} contains the
field \texttt{co\_quorum\_safe}, computed by \texttt{consensus\_obs} from the
underlying consensus state. The proof of admissibility is then immediate:
equality of observations entails equality of the quorum-safety component by
projection. The safety predicate becomes admissible because the observation
carries the safety-relevant distinction directly.

\textit{Restriction.} Rather than asserting global quorum safety, one can
instead assert only what the local observation supports: the node's own quorum
belief, \texttt{local\_quorum\_belief~n}, a projection of \texttt{lv\_quorum}
from the local view record. This predicate is admissible relative to
\texttt{local\_node\_obs~n} because it is a projection of the observation.
The proof requires no axiom. The theorem \texttt{restriction\_is\_genuinely\_weaker}
confirms that local belief and global quorum safety can disagree, so this is a
genuine weakening and not a reformulation.

The theorem \texttt{both\_repairs\_are\_available} states the formal duality
in a single proposition: both \texttt{local\_quorum\_belief~n} and
\texttt{quorum\_safe} are admissible, the former relative to the coarse local
observation, the latter relative to the enriched consensus observation.

% -----------------------------------------------------------------------
\subsection{GPU-Style Memory Observations}

The GPU-style case study concerns the gap between value observations and memory
consistency. A coarse observation may record only the values read and written by
threads. The concrete execution state, however, contains additional ordering and
visibility structure: buffering, cache state, scope, and synchronisation
metadata. A predicate such as sequential consistency of the concurrent execution
is not, in general, determined by the value trace alone.

The admissibility failure has the same form as in the consensus case. Two
executions may have the same observed value trace for a given thread while
disagreeing on whether the execution is sequentially consistent. The observation
map has collapsed a distinction on which the safety predicate depends. The
predicate therefore fails to factor through the thread-local observation.

The GPU case study mechanises both repair paths. Refinement restores
admissibility by enriching the observation with global execution structure.
Restriction keeps the thread-local observation fixed but restricts the domain
to data-race-free executions. On that restricted domain, the
sequential-consistency predicate is stable across observation fibres,
conditional on the DRF-SC bridge stated in \texttt{GPURestriction.v}.

% -----------------------------------------------------------------------
\subsection{Physical Sensing and Boundary Conditions}

The physical-systems case study concerns safety predicates over continuous or
boundary-sensitive states. A controller or monitor rarely observes the full
physical trajectory. It receives sampled, discretised, or otherwise compressed
sensor data. A predicate may require that a variable remain within a safe region
throughout an interval, while the observation records only values at selected
sampling points or coarse boundary summaries.

The resulting admissibility failure is again a failure of factorisation. Two
physical trajectories may induce the same sampled observation while disagreeing
on the safety predicate. One trajectory may remain inside the safe envelope,
while another may violate the envelope between samples or across an unresolved
boundary. The sampled observation is therefore insufficient to determine the
predicate.

This case study mechanises both recovery paths. The coarse sensor observation is
proved inadmissible for the original safety predicate. A restricted predicate
meaningful at the coarse observational grain is then shown to be admissible
relative to that observation. Finally, an enriched sensor observation that adds
the missing physical data restores admissibility for the stronger predicate. The
physical-systems case study therefore demonstrates the full design choice: one
may weaken the guarantee to match the observation, or strengthen the observation
to match the guarantee.

% -----------------------------------------------------------------------
\subsection{Common Structure}

The three case-study families differ in engineering content but not in logical
shape. Consensus hides partition and quorum structure. GPU-style memory hides
ordering and visibility structure. Physical sensing hides inter-sample or
boundary behaviour. In each case, the initial observation map is too coarse for
the safety predicate. The predicate fails to factor because the observation
identifies concrete states that the predicate separates.

This is the common admissibility defect. It is independent of whether the hidden
structure is distributed, concurrent, or physical. The repair is also common:
either choose a predicate that is meaningful at the given observational grain, or
refine the observation until the predicate becomes invariant on observation
fibres. The case studies therefore support the main methodological claim of the
paper. Admissibility is not merely a property of a predicate or a system. It is a
property of the relation between a predicate and an observation interface.

% at the point where the case-study section begins.

\section{Case Studies: Refine or Reject}

The development contains three case-study families: consensus under network
partition, GPU-style relaxed memory, and cyber-physical sensor models. These
examples are deliberately drawn from different engineering domains. Their common
structure is the point. In each case, a safety predicate is naturally formulated
over a concrete state space~\(X\), while the available observation map
\(M : X \to O\) exposes only a coarser interface. The case study then asks
whether the predicate factors through that observation.

The answer is negative for the initial coarse observation in each family. The
coarse observation identifies concrete states that disagree on the relevant
safety predicate. The admissibility failure is therefore structural rather than
accidental. Recovery takes one of two forms: replace the predicate by one that
is meaningful at the observational grain, or refine the observation so that the
safety-relevant distinction is no longer collapsed.

\begin{table}[t]
\centering
\caption{Case-study structure: factorisation failure and admissibility recovery.}
\label{tab:case-study-factorisation}
\small
\begin{tabular}{p{1.8cm}p{2.8cm}p{3.5cm}p{2.9cm}}
\hline
Case study & Hidden distinction & Coarse failure & Recovery \\
\hline
Consensus
  & Partition and quorum structure
  & Local observations do not determine global quorum safety
  & Restrict to local quorum belief; or enrich with quorum-safety data \\[4pt]
GPU memory
  & Buffer, cache, and ordering structure
  & Value traces do not determine consistency behaviour
  & Restrict to DRF executions; or enrich with global ordering metadata \\[4pt]
Physical systems
  & Inter-sample and boundary behaviour
  & Sampled values do not determine continuous safety
  & Use robust observational margins; or enrich sensing and boundary data \\
\hline
\end{tabular}
\end{table}

\paragraph{Mechanisation status.}
The three case-study families instantiate the same factorisation pattern.
The consensus case study mechanises both repair paths: the refinement
path restores admissibility for \(\texttt{quorum\_safe}\) by enriching
the observation (\texttt{ConsensusAdmissibility.v}), and the restriction
path yields an admissible local-quorum-belief predicate at the coarse
observational grain (\texttt{ConsensusRestriction.v}). The GPU case study
mechanises both repair paths: refinement by enriched global execution
observation (\texttt{GPUAdmissibility.v}), and restriction to the subtype
of data-race-free executions (\texttt{GPURestriction.v}). The
physical-systems case study mechanises both paths: a restricted predicate
is admissible at the coarse sensor grain, and an enriched sensor observation
restores admissibility for the stronger predicate. The development
mechanises predicate restriction in all three domains and observational
refinement in all three.

% -----------------------------------------------------------------------
\subsection{Consensus Under Partition}

The consensus case study concerns quorum safety under network partition. The
concrete state records enough information to determine whether the relevant
quorum-safety predicate holds. A coarse observation, however, may expose only
local status information. Such an observation does not determine global quorum
safety. Two concrete states may agree on the local observation while disagreeing
on whether a quorum is safely established.

This failure is not a defect in the proof script. It is the phenomenon the
framework is intended to detect. A uniformity principle for reachable alive
nodes fails in the partition model: different nodes may have different
reachable alive counts. The admissibility proof cannot be closed by that
route.

The mechanised consensus case study proves both repair paths.

\textit{Refinement.} The observation record \texttt{ConsensusObs} contains the
field \texttt{co\_quorum\_safe}, computed by \texttt{consensus\_obs} from the
underlying consensus state. The proof of admissibility is then immediate:
equality of observations entails equality of the quorum-safety component by
projection. The safety predicate becomes admissible because the observation
carries the safety-relevant distinction directly.

\textit{Restriction.} Rather than asserting global quorum safety, one can
instead assert only what the local observation supports: the node's own quorum
belief, \texttt{local\_quorum\_belief~n}, a projection of \texttt{lv\_quorum}
from the local view record. This predicate is admissible relative to
\texttt{local\_node\_obs~n} because it is a projection of the observation.
The proof requires no axiom. The theorem \texttt{restriction\_is\_genuinely\_weaker}
confirms that local belief and global quorum safety can disagree, so this is a
genuine weakening and not a reformulation.

The theorem \texttt{both\_repairs\_are\_available} states the formal duality
in a single proposition: both \texttt{local\_quorum\_belief~n} and
\texttt{quorum\_safe} are admissible, the former relative to the coarse local
observation, the latter relative to the enriched consensus observation.

% -----------------------------------------------------------------------
\subsection{GPU-Style Memory Observations}

The GPU-style case study concerns the gap between value observations and memory
consistency. A coarse observation may record only the values read and written by
threads. The concrete execution state, however, contains additional ordering and
visibility structure: buffering, cache state, scope, and synchronisation
metadata. A predicate such as sequential consistency of the concurrent execution
is not, in general, determined by the value trace alone.

The admissibility failure has the same form as in the consensus case. Two
executions may have the same observed value trace for a given thread while
disagreeing on whether the execution is sequentially consistent. The observation
map has collapsed a distinction on which the safety predicate depends. The
predicate therefore fails to factor through the thread-local observation.

The GPU case study mechanises both repair paths. Refinement restores
admissibility by enriching the observation with global execution structure.
Restriction keeps the thread-local observation fixed but restricts the domain
to data-race-free executions. On that restricted domain, the
sequential-consistency predicate is stable across observation fibres,
conditional on the DRF-SC bridge stated in \texttt{GPURestriction.v}.

% -----------------------------------------------------------------------
\subsection{Physical Sensing and Boundary Conditions}

The physical-systems case study concerns safety predicates over continuous or
boundary-sensitive states. A controller or monitor rarely observes the full
physical trajectory. It receives sampled, discretised, or otherwise compressed
sensor data. A predicate may require that a variable remain within a safe region
throughout an interval, while the observation records only values at selected
sampling points or coarse boundary summaries.

The resulting admissibility failure is again a failure of factorisation. Two
physical trajectories may induce the same sampled observation while disagreeing
on the safety predicate. One trajectory may remain inside the safe envelope,
while another may violate the envelope between samples or across an unresolved
boundary. The sampled observation is therefore insufficient to determine the
predicate.

This case study mechanises both recovery paths. The coarse sensor observation is
proved inadmissible for the original safety predicate. A restricted predicate
meaningful at the coarse observational grain is then shown to be admissible
relative to that observation. Finally, an enriched sensor observation that adds
the missing physical data restores admissibility for the stronger predicate. The
physical-systems case study therefore demonstrates the full design choice: one
may weaken the guarantee to match the observation, or strengthen the observation
to match the guarantee.

% -----------------------------------------------------------------------
\subsection{Common Structure}

The three case-study families differ in engineering content but not in logical
shape. Consensus hides partition and quorum structure. GPU-style memory hides
ordering and visibility structure. Physical sensing hides inter-sample or
boundary behaviour. In each case, the initial observation map is too coarse for
the safety predicate. The predicate fails to factor because the observation
identifies concrete states that the predicate separates.

This is the common admissibility defect. It is independent of whether the hidden
structure is distributed, concurrent, or physical. The repair is also common:
either choose a predicate that is meaningful at the given observational grain, or
refine the observation until the predicate becomes invariant on observation
fibres. The case studies therefore support the main methodological claim of the
paper. Admissibility is not merely a property of a predicate or a system. It is a
property of the relation between a predicate and an observation interface.


\subsection{Comparing the burden across domains}
\label{subsec:domain-burden}

The case studies share a formal pattern. They do not share a mathematical
substrate. In the consensus example, the relevant collapse is induced by local
views under partition. In the GPU example, the relevant collapse is induced by
thread-local observation of an execution whose global consistency status may
differ. In the physical-systems example, the collapse is mediated by sampling,
boundary behaviour, and mereotopological assumptions about regions.

The physical case is therefore the heaviest case, not because the criterion
changes, but because the observation boundary is richer. Spatial connection,
proper containment, inter-sample behaviour, and robust margins must be stated
before the factorisation question is even well posed. This is why the physical
case has explicit RCC-style assumptions and application-level coherence
constraints. The discrete cases need less geometric structure. They are not
therefore more fundamental. They are simpler tests of the same criterion.

This also explains the asymmetry in recovery. Refinement is the natural repair
in the consensus and GPU examples because the missing distinction is a richer
execution or global-state distinction. In the physical example, restriction is
also natural because safe operation is often stated by margins, envelopes, and
boundary conditions. Both repairs are used. They are not used in the same way.

% -----------------------------------------------------------------------
\section{Related Work}
\label{sec:related}
% -----------------------------------------------------------------------

\subsection{Abstraction and abstract interpretation}

The closest broad setting is abstract interpretation. Cousot and Cousot define
program analysis by interpreting concrete computations in an abstract domain,
with soundness controlled by the relation between concrete and abstract
semantics~\cite{cousot1977abstract}. The present paper is narrower. It isolates
one condition inside such abstraction practice: a predicate must be constant on
the observational fibres if it is to be asserted through the observation.

Counterexample-guided abstraction refinement gives a standard method for
repairing abstractions that admit spurious counterexamples~\cite{clarke2000cegar,clarke2003cegar}.
The present criterion can be read as a local algebraic test for when refinement
is necessary. The emphasis is different. CEGAR asks how to refine after an
abstract counterexample. Structural admissibility asks whether the predicate was
observable through the interface in the first place.

\subsection{Runtime verification and monitorability}

Runtime verification studies what can be judged from finite or partial runtime
observations. Bauer, Leucker, and Schallhart give a three-valued account for LTL
and TLTL monitoring, including inconclusive outcomes over partial traces~\cite{bauer2011runtime}.
This paper does not replace monitorability theory. It states a basic obstruction
that any such theory must respect: a monitor cannot distinguish states that its
observation interface identifies.

Functional safety and risk-management practice already distinguish between
changing the operating claim and changing the instrumentation. Standards such
as ISO~31000 and IEC~61508 are not factorisation theories, but they make clear
that risk treatment, safety functions, monitoring, and operating conditions must
be stated together~\cite{iso31000,iec61508}. This paper supplies a small formal
criterion for one version of that practical discipline.

\subsection{Consensus and weak memory}

The consensus case study is intentionally small. It should be read against the
larger literature on distributed agreement, including Lamport's presentation of
Paxos~\cite{lamport2001paxos}. The point is not to mechanise Paxos. The point is
to show how a local observation can collapse states that differ with respect to
a global safety predicate.

The GPU-style case is similarly modest. It is motivated by the general problem
that weak-memory behaviours depend on global consistency constraints not
visible in a local thread trace. Work on x86-TSO and weak memory models gives
the broader context for this issue~\cite{owens2009x86tso,alglave2014herding}.
Here again, the paper does not claim to mechanise a production memory model. It
uses the setting to expose factorisation failure.

\subsection{Type-theoretic and constructive verification}

The mechanisation is carried out in Rocq, formerly Coq. Proof assistants make
assumptions, definitions, and proof obligations explicit. That explicitness is
part of the paper's method. A hidden observational assumption is not harmless
because it permits a theorem to compile. It changes the observation through
which the theorem is being claimed~\cite{coq2024}.

\subsection{Mereotopology and region-based spatial reasoning}

The physical-systems case uses region-based reasoning in the spirit of the
Region Connection Calculus. Randell, Cui, and Cohn's RCC framework treats
spatial reasoning in terms of regions and connection rather than point-set
coordinates~\cite{randell1992rcc}. The present paper uses only a small part of this
tradition. Its role is to make boundary-sensitive observation explicit.

% -----------------------------------------------------------------------
\section{Limitations}
\label{sec:limitations}
% -----------------------------------------------------------------------

The case studies are intentionally minimal. They are designed to isolate
factorisation failure, not to provide full mechanisations of Paxos, a
production GPU memory model, or a continuous control system. This should be
stated plainly. The paper gains nothing by pretending otherwise.

The mechanised recovery paths are asymmetric. Consensus and GPU memory
formalise observational refinement. The physical-systems case formalises both
robust restriction and refinement. This asymmetry is a limitation of coverage,
but it is not a contradiction. It reflects the different repair costs of the
three domains.

The RCC axioms and application-level coherence constraints are explicit
mathematical assumptions. They are not derived from a minimal connection theory.
The paper gives a structural criterion, not an automated refinement algorithm
and not a decision procedure for admissibility.

Finally, warrant debt is not proposed as a quantitative risk metric here. The
paper uses it as a structural name for a failed relation between claim and
observation. Quantitative measures of severity, distance to repair, or expected
runtime harm are future work.

% -----------------------------------------------------------------------
\section{Discussion}
\label{sec:discussion}
% -----------------------------------------------------------------------

\subsection{Observation interfaces as verification boundaries}

A safety property is not only a property of a system. It is also a property of
the relation between the system, the predicate, and the observation interface.
This is the paper's central claim. It is deliberately conservative. If the
observation cannot distinguish two states, then no claim that separates those
states is justified through that observation.

\subsection{Refinement as factorisation repair}

Refinement is often described as the removal of spurious behaviour from an
abstraction. That description is correct, but incomplete for the present
purpose. Here refinement repairs a failed factorisation. It changes the
observation boundary so that a predicate previously invisible at the interface
becomes visible.

This way of speaking also clarifies why refinement is not always the right
repair. Sometimes the better engineering decision is to restrict the claim.
Operating envelopes, robust margins, lockout conditions, and restricted modes
are not second-class repairs. They are ordinary safety devices. What matters is
that the restricted claim actually factors through the observation being used.

\subsection{Warrant debt}

Warrant debt is the operational residue of claiming more than the observation
can support. The term is useful only if it prevents equivocation. If the missing
distinction is irrelevant to the predicate, there is no debt. If the predicate
requires the missing distinction, then calling the discrepancy noise is unsafe.
The proof obligation is to show factorisation, restrict the claim, or refine the
observation.

\paragraph{Future work.}
Future work will investigate quantitative measures of admissibility failure,
including information-theoretic or metric estimates of how much refinement is
needed to separate safe from unsafe fibres of an observation map. Another
natural direction is to mechanise more restriction examples in the discrete case
studies, not because the framework requires symmetry, but because reviewers may
reasonably ask to see both repair forms exercised across more than one domain.

% -----------------------------------------------------------------------
\section{Conclusion}
\label{sec:conclusion}
% -----------------------------------------------------------------------

The paper makes a small claim with sharp consequences. A safety predicate can be
asserted through an observation only if it factors through that observation. If
two concrete states have the same observation but different predicate values,
the observation does not support the claim. Post-processing cannot repair that
loss.

This situation occurs in ordinary safety work. Engineers restrict operating
envelopes. They add sensors. They change thresholds. They refine logs. They
reject claims that the evidence cannot support. The paper does not replace
those practices. It gives one formal reason why they matter.

The warning is simple. An admissibility failure may look like noise. It may look
like harmless underspecification. It may look like a modelling convention. But
if the safety predicate depends on a distinction the observation has collapsed,
the problem is structural. The correct response is not optimism. It is
restriction, rejection, or refinement.

% -----------------------------------------------------------------------
\appendix

\section{Assumptions}
\label{app:assumptions}
% -----------------------------------------------------------------------

The development contains no admitted lemmas, 39 compiled Rocq source files, and
eight named axioms. Every axiom is documented in \texttt{ASSUMPTIONS.md} in the
accompanying repository. The audit commands that verify this status are:

\begin{verbatim}
find . -name '*.v' | wc -l          # 39
find . -name '*.vo' | wc -l         # 39
grep -RniE '\bAdmitted\b' . --include='*.v'   # (no output)
grep -RniE '\bAxiom\b' . --include='*.v'      # 8 matches
\end{verbatim}

Before submission, this appendix should list all eight axiom names directly.
The current paper text gives the role of all eight, but the exact names of the
five application-level axioms must be copied from \texttt{ASSUMPTIONS.md}.

% -----------------------------------------------------------------------
\section{Rocq File Map}
\label{app:filemap}
% -----------------------------------------------------------------------

\begin{table}[h]
\centering
\begin{tabular}{lll}
\toprule
Component & Files & Role \\
\midrule
Foundation & \texttt{Foundation/*.v} & Observation, factorisation, mereotopology \\
Admissibility & \texttt{Admissibility/*.v} & Core admissibility theory \\
Runtime & \texttt{Runtime/*.v} & Monitor states, fatigue, certificates \\
Automation & \texttt{Automation/*.v} & Custom proof tactics \\
Consensus & \texttt{CaseStudies/Consensus/*.v} & Quorum-safety admissibility \\
GPU & \texttt{CaseStudies/GPU/*.v} & Thread-local inadmissibility, global refinement \\
Physical & \texttt{CaseStudies/PhysicalSystems/*.v} & Sensor inadmissibility, restriction, refinement \\
Extraction & \texttt{Extraction/*.v} & Verified OCaml extraction bridge \\
\bottomrule
\end{tabular}
\caption{Rocq source file map.}
\label{tab:filemap}
\end{table}

% -----------------------------------------------------------------------
% Bibliography
% -----------------------------------------------------------------------
\begin{thebibliography}{99}

\bibitem{alglave2014herding}
Jade Alglave, Luc Maranget, and Michael Tautschnig.
\newblock Herding cats: modelling, simulation, testing, and data mining for weak memory.
\newblock \emph{ACM Transactions on Programming Languages and Systems}, 36(2), 2014.

\bibitem{bauer2011runtime}
Andreas Bauer, Martin Leucker, and Christian Schallhart.
\newblock Runtime verification for LTL and TLTL.
\newblock \emph{ACM Transactions on Software Engineering and Methodology}, 20(4):14:1--14:64, 2011.

\bibitem{clarke2000cegar}
Edmund M. Clarke, Orna Grumberg, Somesh Jha, Yuan Lu, and Helmut Veith.
\newblock Counterexample-guided abstraction refinement.
\newblock In \emph{Computer Aided Verification}, LNCS 1855, pages 154--169. Springer, 2000.

\bibitem{clarke2003cegar}
Edmund M. Clarke, Orna Grumberg, Somesh Jha, Yuan Lu, and Helmut Veith.
\newblock Counterexample-guided abstraction refinement for symbolic model checking.
\newblock \emph{Journal of the ACM}, 50(5):752--794, 2003.

\bibitem{cousot1977abstract}
Patrick Cousot and Radhia Cousot.
\newblock Abstract interpretation: a unified lattice model for static analysis of programs by construction or approximation of fixpoints.
\newblock In \emph{Proceedings of the 4th ACM Symposium on Principles of Programming Languages}, pages 238--252, 1977.

\bibitem{iec61508}
International Electrotechnical Commission.
\newblock \emph{IEC 61508: Functional safety of electrical/electronic/programmable electronic safety-related systems}.
\newblock IEC, current multi-part standard.

\bibitem{iso31000}
International Organization for Standardization.
\newblock \emph{ISO 31000:2018 Risk management, Guidelines}.
\newblock ISO, 2018.

\bibitem{lamport2001paxos}
Leslie Lamport.
\newblock Paxos made simple.
\newblock \emph{ACM SIGACT News}, 32(4):51--58, 2001.

\bibitem{owens2009x86tso}
Scott Owens, Susmit Sarkar, and Peter Sewell.
\newblock A better x86 memory model: x86-TSO.
\newblock In \emph{Theorem Proving in Higher Order Logics}, LNCS 5674, pages 391--407. Springer, 2009.

\bibitem{randell1992rcc}
David A. Randell, Zhan Cui, and Anthony G. Cohn.
\newblock A spatial logic based on regions and connection.
\newblock In \emph{Proceedings of the 3rd International Conference on Principles of Knowledge Representation and Reasoning}, pages 165--176, 1992.

\bibitem{coq2024}
The Coq Development Team.
\newblock \emph{The Coq Proof Assistant}.
\newblock Zenodo, 2024.

\end{thebibliography}

\end{document}
